\section{Требования к интерфейсу}
\label{sec:interface_requirements}

В данном разделе описываются требования к пользовательскому интерфейсу разрабатываемой системы, включая его архитектуру, структуру, органы управления и специальные требования, обусловленные особенностями предметной области и целевой аудитории.

\subsection{Тип и платформа интерфейса}

Система реализует \textbf{минималистичный гибридный интерфейс}, состоящий из двух компонентов:

\begin{enumerate}
    \item \textbf{Конфигурационный интерфейс (консольный/оконный)} --- предназначен для первоначальной настройки системы перед началом игры. Реализован в виде командной строки с возможностью расширения до простого окна на базе PyQt6. Обеспечивает:
    \begin{itemize}
        \item выбор покер-рума из списка поддерживаемых;
        \item автоматическое определение окна клиента покера по заголовку;
        \item загрузку пользовательских диапазонов в текстовом формате Flopzilla или выбор предустановленного профиля;
        \item базовую настройку параметров HUD: включение/отключение, позиция на экране, уровень прозрачности.
    \end{itemize}
    
    \item \textbf{Игровой интерфейс (текстовый HUD Overlay)} --- прозрачное оверлейное окно, отображаемое поверх клиента покер-рума в реальном времени. Имеет строго текстовую структуру без графических элементов, что обеспечивает минимальное потребление ресурсов и максимальную скорость обновления. Обеспечивает:
    \begin{itemize}
        \item отображение текущих карт игрока и позиции за столом в текстовом формате;
        \item индикацию статуса парсинга (успешно/ошибка/ожидание);
        \item визуализацию рекомендации по действию;
        \item отображение ключевых параметров стола в компактном текстовом виде.
    \end{itemize}
\end{enumerate}

Оба компонента функционируют в операционной системе Windows 10/11, что обусловлено распространённостью данной платформы среди пользователей онлайн-покера и доступностью необходимых системных вызовов для захвата изображения окон.

\subsection{Специальные требования к интерфейсу}

Требования к интерфейсу сформированы с учётом следующих ограничений и особенностей:

\begin{itemize}
    \item \textbf{Требования к целевой аудитории}: система ориентирована на начинающих игроков, поэтому интерфейс должен быть предельно простым и не требующим обучения. Все элементы отображаются в виде читаемого текста без избыточной графики.
    
    \item \textbf{Ограничения игрового процесса}: во время раздачи игрок принимает решения в условиях ограниченного времени (обычно 15--30 секунд). Следовательно, рекомендация должна отображаться не позднее чем через 500 мс после наступления хода игрока, а время реакции интерфейса на изменение состояния стола не должно превышать 200 мс.
    
    \item \textbf{Технические ограничения}: оверлейное окно должно быть \textit{click-through} (пропускать клики мыши), чтобы не блокировать взаимодействие пользователя с клиентом покера. Прозрачность фона оверлея должна настраиваться в диапазоне 0--90\% для обеспечения читаемости текста при любом оформлении покер-рума.
    
    \item \textbf{Требования к надёжности}: в случае сбоя парсинга интерфейс должен отображать понятное текстовое сообщение об ошибке без прекращения работы системы, позволяя пользователю продолжить игру вручную.
\end{itemize}

\subsection{Структура конфигурационного интерфейса}

Конфигурационный интерфейс реализован в минималистичном стиле и предоставляет пользователю только необходимые элементы управления.

Основные параметры, настраиваемые при запуске системы:

\begin{itemize}
    \item \textbf{Выбор покер-рума}: Параметр задаётся аргументом командной строки \texttt{--room} или через оконный интерфейс;
    
    \item \textbf{Определение окна}: система автоматически находит окно покер-клиента по ключевым словам в заголовке. При необходимости пользователь может указать заголовок вручную или сам клиент;
    
    \item \textbf{Управление диапазонами}: 
    \begin{itemize}
        \item по умолчанию используются предустановленные диапазоны для каждой позиции;
        \item предусмотрена возможность загрузки пользовательского файла с диапазонами в формате, совместимом с Flopzilla (строка вида \texttt{AA-55,AKs-A2s,KQo});
    \end{itemize}
    
    \item \textbf{Настройки HUD}:
    \begin{itemize}
        \item позиция оверлея задаётся координатами или привязкой к углу экрана;
        \item прозрачность фона настраивается значением от 0 до 1 (по умолчанию 0.7).
    \end{itemize}
\end{itemize}

Интерфейс следует принципу \textit{конвенции вместо конфигурации}: большинство параметров имеют разумные значения по умолчанию, что позволяет начать работу с системой без дополнительной настройки.

\subsection{Структура игрового интерфейса (текстовый HUD)}

Оверлейное окно имеет компактную вертикальную текстовую компоновку. Все элементы представлены в виде строк текста с фиксированным шрифтом, что обеспечивает предсказуемое выравнивание и высокую читаемость.

Структура оверлея (сверху вниз):

\begin{enumerate}
    \item \textbf{Строка статуса}:
    \begin{verbatim}
    [OK] PokerOCR v1.0 | ReplayPoker
    \end{verbatim}
    Индикатор в квадратных скобках отражает статус работы системы:
    \begin{itemize}
        \item \texttt{[OK]} --- парсинг работает штатно;
        \item \texttt{[WAIT]} --- ожидание данных (карты не распознаны);
        \item \texttt{[ERR]} --- ошибка парсинга (требуется внимание пользователя).
    \end{itemize}
    
    \item \textbf{Блок информации о герое}:
    \begin{verbatim}
    Hero: Ks Qs (CO) | Bet: 0
    \end{verbatim}
    Формат: \texttt{Hero: <карты> (<позиция>) | Bet: <ставка>}. Карты отображаются в компактном формате (\texttt{Ks} = король пик, \texttt{Qd} = дама бубен).
    
    \item \textbf{Блок параметров стола}:
    \begin{verbatim}
    Blinds: 100/200 | Players: 4 | Raises: 1
    \end{verbatim}
    Краткая сводка: размер блайндов, количество активных игроков, число рейзов в текущей раздаче.
    
    \item \textbf{Блок рекомендации} (ключевой элемент, выделен визуально):
    \begin{verbatim}
    >>> RAISE [HIGH] KQs in CO range
    \end{verbatim}
    Формат: \texttt{>>> <действие> [<уверенность>] <обоснование>}.
    \begin{itemize}
        \item Действие: \texttt{FOLD}, \texttt{CALL}, \texttt{RAISE};
        \item Уверенность: \texttt{[LOW]}, \texttt{[MED]}, \texttt{[HIGH]};
        \item Обоснование: краткий текст, объясняющий рекомендацию.
    \end{itemize}
    
    \item \textbf{Разделитель}: строка из символов \texttt{=}, визуально отделяющая рекомендацию от технической информации.
\end{enumerate}

Пример полного вывода оверлея:
\begin{verbatim}
[OK] PokerOCR v1.0 | ReplayPoker
Hero: Ks Qs (CO) | Bet: 0
Blinds: 100/200 | Players: 4 | Raises: 1
>>> RAISE [HIGH] KQs in CO range
==================================================
\end{verbatim}

Все текстовые элементы используют моноширинный шрифт (например, \texttt{Consolas} или \texttt{Courier New}) для обеспечения точного выравнивания. Цветовое оформление:
\begin{itemize}
    \item основной текст --- белый или светло-серый;
    \item рекомендация \texttt{RAISE} --- зелёный;
    \item рекомендация \texttt{FOLD} --- красный;
    \item рекомендация \texttt{CALL} --- жёлтый;
    \item статус ошибки --- красный мигающий.
\end{itemize}

\subsection{Динамика работы интерфейса}

Работа интерфейса описывается конечным автоматом с тремя основными состояниями.

\textbf{Состояние 1: Инициализация}
\begin{itemize}
    \item Система загружает конфигурацию, инициализирует OCR-модель, находит окно покера.
    \item В оверлее отображается: \texttt{[INIT] Loading...}
    \item При успехе --- переход в состояние \textit{Ожидание}.
    \item При ошибке --- повтор попытки или отображение \texttt{[ERR] Window not found}.
\end{itemize}

\textbf{Состояние 2: Ожидание}
\begin{itemize}
    \item Система в реальном времени обрабатывает кадры, но ход ещё не наступил.
    \item В оверлее отображается актуальная информация о картах и столе, но блок рекомендации пуст или содержит \texttt{>>> -- waiting --}.
    \item При детекции наступления хода игрока (появление кнопок действия) --- переход в состояние \textit{Рекомендация}.
\end{itemize}

\textbf{Состояние 3: Рекомендация}
\begin{itemize}
    \item Система анализирует ситуацию, сопоставляет руку с диапазоном, формирует совет.
    \item В оверлее отображается полноценная рекомендация с действием, уверенностью и обоснованием.
    \item После совершения игроком действия или по таймауту (30 с) --- возврат в состояние \textit{Ожидание}.
\end{itemize}

Переходы между состояниями происходят автоматически на основе данных парсинга и не требуют вмешательства пользователя.

\subsection{Требования к доступности и эргономике}

Для обеспечения удобства использования системой предусмотрены следующие решения:

\begin{itemize}
    \item \textbf{Минимизация визуального шума}: интерфейс содержит только необходимую информацию, без декоративных элементов, анимаций или избыточной графики. Это снижает когнитивную нагрузку на игрока в условиях ограниченного времени.
    
    \item \textbf{Читаемость текста}: используется моноширинный шрифт размером не менее 10 pt, высокий контраст между текстом и фоном, возможность настройки прозрачности фона оверлея.
    
    \item \textbf{Позиционирование}: оверлей по умолчанию размещается в правом верхнем углу экрана, где обычно не располагаются ключевые элементы покер-клиента (карты, кнопки действий). Пользователь может изменить позицию через настройки.
    
\end{itemize}

Таким образом, разработанный интерфейс обеспечивает баланс между функциональностью, скоростью реакции и минимализмом, что соответствует требованиям к вспомогательному программному обеспечению для онлайн-покера.
